%=============================================================================
% Section 2: Formal Architectural Model and Problem Formulation
%=============================================================================
\section{Formal Architectural Model and Problem Formulation}

\subsection{Logical Model: $S = (A, P, E)$}

The AI execution architecture is defined as the following logical tuple:
\[
S = (A, P, E)
\]
where:
\begin{itemize}
    \item $A$ represents the \textbf{Agent Intelligence Layer} (reasoning and planning),
    \item $P$ represents the \textbf{Contact Layer} (routing and decision-making),
    \item $E$ represents the \textbf{Execution Layer} (capability invocation).
\end{itemize}

The model separates the AI system into three distinct layers, each with clearly defined responsibilities:

\textbf{Agent Layer:} Agents represent intelligent behavior. They reason about tasks, plan actions, and decide which capabilities should be invoked. This layer handles:
\begin{itemize}
    \item Understanding user goals
    \item Planning sequences of actions
    \item Making decisions about capability requirements
    \item Managing conversation context and memory
\end{itemize}

\textbf{Contact Layer ($P$):} The Contact Layer determines \emph{how} capabilities are selected and composed. This layer is responsible for:
\begin{itemize}
    \item Capability selection based on agent intent
    \item Provider preference and selection
    \item Routing strategies
    \item Cost or latency optimization
    \item Policy enforcement
\end{itemize}

\textbf{Execution Layer:} The execution layer is responsible \emph{only} for executing capability calls. It does not implement routing policies, provider selection, or agent logic. This design keeps the runtime minimal and reusable.

\subsection{Logical Model vs.\ Implementation Stack}

In this paper, $S = (A, P, E)$ is the \textbf{logical model}. The five-layer AI Native Stack presented later is an \textbf{implementation stack} used to realize this model in production systems.
\begin{itemize}
    \item The logical model specifies invariants and boundaries.
    \item The implementation stack specifies deployment surfaces and engineering composition.
\end{itemize}

This distinction avoids a common modeling conflict: a system can be five-layer in implementation while remaining three-layer in execution logic.

\subsection{Rationale for Separation}

This separation addresses a critical architectural problem observed in existing frameworks. Many systems conflate orchestration responsibilities with execution mechanisms, resulting in:
\begin{itemize}
    \item \textbf{Framework Lock-in:} Applications become tightly coupled to specific orchestration frameworks
    \item \textbf{Limited Reusability:} Runtime components cannot be easily reused across different applications
    \item \textbf{Complexity Creep:} Execution runtimes grow to include policy logic, making them difficult to maintain
\end{itemize}

By separating these concerns, we enable:
\begin{itemize}
    \item \textbf{Runtime Portability:} Execution runtimes can be reused across different orchestration strategies
    \item \textbf{Policy Flexibility:} Organizations can implement custom routing policies without modifying the runtime
    \item \textbf{Clear Boundaries:} Each component has well-defined inputs, outputs, and responsibilities
\end{itemize}
